\section{The Settlement Layer Interface}
\label{sec:sec12}
\label{sec:settlement}
%==============================================================================

The Ledger records economic events; the settlement layer executes real-world transfers. The boundary between them is a \emph{projection}: a pure, total function that reads one committed transaction and produces an optional settlement instruction. Everything the settlement layer needs is already present in the transaction; everything the settlement layer adds---custodian accounts, CSD connectivity, message formatting---lies outside the Ledger's scope.

\subsection{The Settlement Projection}
\label{sec:settlement-projection}

\begin{definition}[Settlement Projection]
\label{def:settle-projection}
The settlement projection
\[
\texttt{settleProjection} : \texttt{SettlementTx} \to \texttt{Maybe}\ \texttt{SettlementInstruction}
\]
is pure, deterministic, stateless, and idempotent. It is total: defined on every settlement transaction, returning \texttt{Just} an instruction for the settlable classes (\texttt{SETTLEMENT}, \texttt{COLLATERAL}) and \texttt{Nothing} for the rest.
\end{definition}

Its input, \texttt{SettlementTx}, is a \emph{boundary type}, not the core transaction of Section~\ref{sec:ledger}. The core \texttt{Transaction} is the one atomic event recorded in the log---an ordered list of moves over ledger units, with the state delta they ride. \texttt{SettlementTx} is the settlement-layer view onto which a committed core transaction \emph{projects}: the transaction id, its settlement class, and its moves with each move's unit already resolved to its settlement identity (\texttt{Asset}---an ISIN or a currency). The projection from core \texttt{Transaction} to \texttt{SettlementTx} is the unit-to-asset resolution the contract performs at move generation; \texttt{settleProjection} then carries \texttt{SettlementTx} to the wire instruction. The two never share a name: there is one \texttt{Transaction} (the log event), and \texttt{SettlementTx} is what the boundary reads of it.

The projection reads only from the transaction:
\begin{itemize}[nosep]
\item \textbf{Transaction id} --- the settlement reference (\texttt{EndToEndId} in ISO~20022).
\item \textbf{Each move's source wallet, destination wallet, settlement asset, and quantity} --- the economic content of the transfer.
\item \textbf{The CDM event payload} (Section~\ref{sec:cdm}) --- settlement date, counterparty identifiers (LEI, BIC), and execution venue (MIC).
\end{itemize}

It reads nothing else: not current wallet balances, not other transactions, not any external system. Purity follows: the same settlement transaction always yields the same instruction, the function is safe to re-run after a failure, and the settlement layer need not understand conservation, unit state, or lifecycle logic.

The projection, types first. A settlable move's \texttt{smAsset} is its unit's settlement identity---the security ISIN or the currency the unit denotes, resolved from the unit's registered terms (Section~\ref{sec:sec04}) and carried into the \texttt{SettlementTx} by the contract---so the projection needs no off-transaction lookup. The settlement \emph{type} and the \emph{legs} are one value: each constructor of \texttt{SettlementLegs} carries exactly the legs its type names, so an instruction with no legs cannot be written and the type can never disagree with the legs.

\begin{lstlisting}[language=Haskell]
-- The settlement classification of a unit: a transferable security, or a
-- currency. Every settlable move is exactly one or the other. (CashCcy, not
-- Cash, so the settlement currency case never collides with a scalar.)
data Asset = Security ISIN | CashCcy Currency

-- A settlement-layer move: it ranges over the settlement identity (Asset),
-- not the core ledger UnitId. It is what a core Move projects to once the
-- contract has resolved the unit to its settlement asset.
data SettleMove = SettleMove { smFrom, smTo :: WalletId, smAsset :: Asset, smQty :: Qty }

-- Settlability is a field set by the contract, never inferred. Total over all
-- five classes; only SETTLEMENT and COLLATERAL settle. LIFECYCLE/CORRECTION
-- "sometimes" settle only by DECOMPOSING into separate transactions at the source.
data TxClass = SETTLEMENT | COLLATERAL | LIFECYCLE | ACCOUNTING | CORRECTION
settles :: TxClass -> Bool
settles SETTLEMENT = True
settles COLLATERAL = True
settles LIFECYCLE  = False
settles ACCOUNTING = False
settles CORRECTION = False

data CdmPayload   = CdmPayload   { cdmTradeDate, cdmSettleDate :: Day
                                , cdmCounterparty :: LEI, cdmVenue :: MIC }

-- The boundary type. NOT the core Transaction (the log event of Sec.~ledger):
-- it is the settlement-layer view a committed core Transaction PROJECTS onto,
-- its moves already resolved to settlement Assets.
data SettlementTx = SettlementTx { stId :: TxId, stClass :: TxClass
                                 , stMoves :: [SettleMove], stCdm :: CdmPayload }

data SecuritiesLeg = SecuritiesLeg { slIsin :: ISIN, slQty :: Qty
                                   , slDeliver, slReceive :: WalletId }
data CashLeg       = CashLeg       { clCcy :: Currency, clAmount :: Qty
                                   , clPayer, clReceiver :: WalletId }

-- type and legs are ONE value: "an instruction with no legs" is unrepresentable.
data SettlementLegs = DvP SecuritiesLeg CashLeg | FoP SecuritiesLeg | CashOnly CashLeg

data SettlementType = DVP | FOP | CASH
settlementType :: SettlementLegs -> SettlementType   -- a projection, not a stored field
settlementType (DvP _ _)    = DVP
settlementType (FoP _)      = FOP
settlementType (CashOnly _) = CASH

data SettlementInstruction = SettlementInstruction      -- plain data, no behaviour
  { siRef :: TxId, siTradeDate, siSettleDate :: Day
  , siLegs :: SettlementLegs, siCounterparty :: LEI, siVenue :: MIC }

-- A security move -> securities leg (source delivers, destination receives);
-- a cash move -> cash leg (source pays, destination receives). Total over Asset.
classify :: SettleMove -> Either SecuritiesLeg CashLeg
classify (SettleMove from to (Security i) q) = Left  (SecuritiesLeg i q from to)
classify (SettleMove from to (CashCcy c)  q) = Right (CashLeg       c q from to)

-- Nothing when there is no standard leg shape: no moves (a contract error), or a
-- non-decomposed multi-leg structure. Both give no instruction, never a partial one.
legsOf :: [SettleMove] -> Maybe SettlementLegs
legsOf ms = case partitionEithers (map classify ms) of
  ([s],[c]) -> Just (DvP s c)
  ([s],[])  -> Just (FoP s)
  ([],[c])  -> Just (CashOnly c)
  _         -> Nothing

settleProjection :: SettlementTx -> Maybe SettlementInstruction   -- pure, total
settleProjection tx
  | settles (stClass tx) = mk <$> legsOf (stMoves tx)
  | otherwise            = Nothing
  where mk legs = SettlementInstruction (stId tx) (cdmTradeDate (stCdm tx))
          (cdmSettleDate (stCdm tx)) legs (cdmCounterparty (stCdm tx)) (cdmVenue (stCdm tx))
\end{lstlisting}

\begin{lstlisting}[language=Haskell]
-- Behaviour, made concrete  (settlementType . siLegs <$> settleProjection tx):
--   security + cash move  (a DvP equity trade) -> Just DVP
--   a securities move alone                     -> Just FOP
--   a cash move alone     (a margin call)       -> Just CASH
--   an ACCOUNTING transaction                   -> Nothing   -- not settlable
--   a SETTLEMENT with no moves                  -> Nothing   -- contract error, no instruction
\end{lstlisting}

\begin{remark}
The projection is total and deterministic by construction. \texttt{settles} classifies every \texttt{TxClass}; \texttt{legsOf} is exhaustive, its wildcard absorbing both the empty and the non-decomposed multi-leg cases as \texttt{Nothing}; and \texttt{settleProjection} reads only \texttt{tx}---no clock, no balances, no external system---so the same settlement transaction always yields the same instruction and re-running after a failure is safe.
\end{remark}

\begin{definition}[SettlementInstruction]
\label{def:settlement-instruction}
The \texttt{SettlementInstruction} is plain data with no behaviour (the \texttt{data} declaration above): instruction id, trade and settlement dates, the settlement legs, counterparty LEI, and execution venue MIC. The settlement type $\in \{\texttt{DVP}, \texttt{FOP}, \texttt{CASH}\}$ is a \emph{projection} of which legs are present (\texttt{settlementType}), never an independent field that could disagree with them. Both legs present give \texttt{DVP}; securities alone give \texttt{FOP} (Free of Payment); cash alone gives \texttt{CASH}; no legs is unrepresentable, so a settlable transaction with no moves yields \texttt{Nothing}---a contract error, not an instruction.
\end{definition}

The standard cases---DvP equity trades, FOP transfers, cash-only settlements---carry at most one securities leg and one cash leg. Multi-leg structures decompose at the source, not in the projection: the smart contract emits one \texttt{SETTLEMENT} transaction per independently settleable leg. A repo emits two---an opening leg and a closing leg---each a standard DvP structure.

\subsection{Boundary Separation}
\label{sec:settlement-boundary}

\begin{principle}[Boundary Separation]
\label{prin:settlement-boundary}
The Ledger provides \emph{what} settles: parties, securities, quantities, dates. The settlement layer provides \emph{how}: standing settlement instructions (SSIs), custodian account identifiers, CSD connectivity, ISO~20022 message generation, and exception management. Neither side reads the other's internals.
\end{principle}

The \texttt{SettlementInstruction} is the sole object shared across the boundary. Wire-message generation is two steps: projection, then enrichment. The projection produces the instruction from the transaction; the settlement layer enriches it with the operational detail the Ledger does not hold---SSIs, CSD participant ids, settlement priority---before generating the ISO~20022 message (e.g.\ \texttt{sese.023}). This is the precise sense of ``message generation is projection, not translation'' (Section~\ref{sec:implementation}): the projection step is the Ledger's; enrichment is the settlement layer's.

\subsection{Transaction Types and Settlability}
\label{sec:settlement-types}

Settlability is a first-class field on the transaction, assigned by the generating smart contract at the point of move generation, never inferred at runtime---the contract knows the business intent.

\begin{center}
\begin{tabular}{l l p{7.5cm}}
\toprule
\textbf{Type} & \textbf{Settles?} & \textbf{Description} \\
\midrule
\texttt{SETTLEMENT} & Yes & Maps to a real-world settlement instruction (e.g.\ equity trade $\to$ \texttt{sese.023}). \\
\texttt{COLLATERAL} & Yes & Settles through the collateral channel (e.g.\ margin call $\to$ cash transfer). \\
\texttt{LIFECYCLE} & Sometimes & Decomposes into settlement and accounting sub-transactions. \\
\texttt{ACCOUNTING} & No & Internal adjustment, no external transfer (e.g.\ PnL crystallisation, revaluation). \\
\texttt{CORRECTION} & Sometimes & Compensating transaction; settles a reversal only if the original already settled. \\
\bottomrule
\end{tabular}
\end{center}

The projection returns \texttt{Just} for \texttt{SETTLEMENT} and \texttt{COLLATERAL}, \texttt{Nothing} for the rest. When a \texttt{LIFECYCLE} event requires both settlement and accounting, the smart contract emits both transactions: a \texttt{SETTLEMENT} for the legs that transfer externally, an \texttt{ACCOUNTING} for the internal adjustment. Physical exercise of an option, for instance, emits a \texttt{SETTLEMENT} (deliver shares, receive cash) and a separate \texttt{ACCOUNTING} (close the option position, crystallise PnL). The decomposition is the contract's responsibility, not the projection's.

\subsection{Lifecycle-Originated Settlement}
\label{sec:settlement-lifecycle}

Some settlements originate from internal contract logic rather than an external instruction: an option exercise, a coupon falling due, a barrier breach. The Ledger is the originator, so there is no external instruction to match. The projection treats these identically---the transaction is tagged \texttt{SETTLEMENT} by the contract, the projection extracts the instruction, the settlement layer enriches and delivers it. The sole difference is provenance: execution-originated instructions match an external trade confirmation; lifecycle-originated instructions are generated unilaterally. The settlement layer receives the instruction with no awareness of its origin.

One lifecycle transaction may yield several instructions where it carries several legs. Physical exercise of an equity option with lot-size constraints (Section~\ref{sec:contracts}) carries a securities delivery (whole lots) and a cash payment (strike value plus any fractional residual); the projection classifies each move by unit type and emits the matching instruction for each.

\subsection{DvP: The Two-Level Guarantee}
\label{sec:settlement-dvp}

Delivery versus Payment (DvP) is the principle that securities delivery and cash payment occur together or not at all. The framework provides it at two independent, mutually reinforcing levels.

\paragraph{Ledger-level.} A DvP trade is one transaction holding both the securities move and the cash move. The executor commits all moves atomically---both take effect or neither does (Section~\ref{sec:ledger}). This is a structural consequence of the transaction primitive, not a runtime check; the smart-contract specification (Section~\ref{sec:contracts}) supplies the other half by requiring both legs within one transaction.

\paragraph{Settlement-level.} The CSD or CCP provides DvP at the real-world level: securities delivery and cash payment settle simultaneously on its books. This depends on external infrastructure and lies outside the Ledger's control.

\paragraph{The temporal gap.} Between trade date and settlement date the Ledger shows the economic position under trade-date accounting (Section~\ref{sec:intro}); the CSD confirms the real-world transfer at settlement. The gap is not a DvP failure but the normal state under trade-date accounting, where exposure exists from execution.

\paragraph{Settlement failure.} On failure (the counterparty cannot deliver), the Ledger records a settlement-failure event and does \emph{not} reverse the economic position: the exposure existed from trade date regardless of settlement outcome, and persists until the failure resolves through buy-in, cancellation, or partial delivery. The discrepancy between Ledger and custodian positions is detected by virtual-wallet comparison (Section~\ref{sec:implementation}).

\paragraph{External regimes.} Legal title, beneficial ownership, fail treatment, and the regulatory-capital treatment of unsettled transactions are governed by external frameworks, not by the Ledger. The Ledger supplies the position and timing data those treatments require; it does not determine the applicable regime.

\begin{remark}
The mechanics of the trade-date-to-settlement-date interval---deferred settlement, fail lifecycle, buy-in, partial delivery, and re-instruction---are specified in the standalone deferred-settlement specification. This section fixes only the interface: what the projection produces and how settlement status returns. The satellite is out of scope here.
\end{remark}

\subsection{Gross-to-Net Reconciliation}
\label{sec:settlement-netting}

The Ledger records gross transactions, one per trade. The settlement layer may net: for trades in the same security with the same counterparty on the same settlement date, it produces a single net instruction. Netting is a settlement-layer optimisation; the Ledger's gross records are authoritative and are never modified by it.

Gross and net reconcile by an algebraic identity, not a heuristic. For each (security, counterparty, settlement date) group, the algebraic sum of gross instruction quantities---deliveries positive, receipts negative---equals the net instruction quantity. The sum is a \emph{group homomorphism}---a structure-preserving map, here meaning summing the gross legs gives the net leg---written \texttt{foldMap} from the gross legs into the additive group of quantities. A violation localises the fault: either the projection produced an incorrect instruction (a Ledger fault) or the netting computation is wrong (a settlement-layer fault).

\begin{lstlisting}[language=Haskell]
-- Signed securities quantity from one party's view: delivery positive, receipt negative.
signedFor :: WalletId -> SecuritiesLeg -> Qty
signedFor me l | slDeliver l == me = slQty l
               | slReceive l == me = negQty (slQty l)
               | otherwise         = mempty

-- The net leg's signed quantity equals the sum of the gross legs'. An identity,
-- checked, not a reconciliation heuristic.
reconciles :: WalletId -> [SecuritiesLeg] -> SecuritiesLeg -> Bool
reconciles me gross net = foldMap (signedFor me) gross == signedFor me net
\end{lstlisting}

\subsection{The Confirmation Return Path}
\label{sec:settlement-confirmation}

Settlement confirmations (\texttt{sese.025} for securities, \texttt{camt.054} for cash) flow back into the Ledger as lifecycle events recording the confirmation timestamp and external reference. No moves are generated---positions were already correct from trade date. The settlement status advances:

\begin{center}
\begin{tabular}{l p{10.5cm}}
\toprule
\textbf{Status} & \textbf{Meaning} \\
\midrule
\texttt{EXECUTED}   & Committed in the Ledger; the settlement layer has not yet processed it. \\
\texttt{INSTRUCTED} & The settlement layer has submitted the instruction; the CSD has not confirmed. \\
\texttt{SETTLED}    & The CSD confirms successful settlement; custody now matches Ledger positions. \\
\texttt{FAILED}     & The CSD reports a failure. Ledger positions remain correct (economic exposure); the external transfer has not occurred. A failed trade may re-instruct (return to \texttt{INSTRUCTED} after buy-in or extension), partially settle (a lifecycle event), or cancel (a \texttt{CORRECTION} transaction reversing the original moves). \\
\bottomrule
\end{tabular}
\end{center}

\noindent The transitions are \texttt{EXECUTED} $\to$ \texttt{INSTRUCTED} $\to$ \texttt{SETTLED}, with \texttt{INSTRUCTED} $\to$ \texttt{FAILED} as the alternate terminal. Each transition is itself a logged lifecycle event in the move stream. Settlement status is therefore a projection of those events, not a separate mutable store: replaying the events to any point rebuilds the status that held then, and the stored value is overwritten in place only as a read cache. This completes the traceability chain of Section~\ref{sec:implementation}: smart contract $\to$ move $\to$ settlement instruction $\to$ external confirmation $\to$ status.

%==============================================================================
